\documentclass[a4paper,12pt]{article}
\usepackage[utf8]{inputenc}
\usepackage[russian]{babel}
\usepackage{amsmath, amssymb}

\title{Задание по циклическим кодам}
\author{}
\date{}

\begin{document}

\maketitle

\section{Факторизация многочлена}

Рассмотрим многочлен $x^7 - 1$ над полем $\mathbb{F}_2$.
Поскольку в этом поле выполняется равенство $-1 = 1$, имеем:
\[
x^7 - 1 = x^7 + 1.
\]

Разложение данного многочлена на неприводимые множители имеет вид:
\[
x^7 + 1 = (x + 1)(x^3 + x + 1)(x^3 + x^2 + 1).
\]

\section{Порождающие полиномы}

Каждый делитель многочлена $x^7 + 1$ задаёт циклический код длины $n=7$.
Исключая тривиальные случаи, получаем следующие варианты порождающих полиномов:

\begin{itemize}
    \item $g_1(x) = x + 1$,
    \item $g_2(x) = x^3 + x + 1$,
    \item $g_3(x) = x^3 + x^2 + 1$,
    \item $g_4(x) = (x + 1)(x^3 + x + 1) = x^4 + x^3 + x^2 + 1$,
    \item $g_5(x) = (x + 1)(x^3 + x^2 + 1) = x^4 + x^2 + x + 1$,
    \item $g_6(x) = (x^3 + x + 1)(x^3 + x^2 + 1)$,
\end{itemize}

где последний полином можно записать как:
\[
g_6(x) = x^6 + x^5 + x^4 + x^3 + x^2 + x + 1.
\]

\section{Параметры кодов}

Для каждого порождающего полинома можно определить параметры кода.
Размерность вычисляется по формуле:
\[
k = n - \deg g(x).
\]

\begin{center}
\begin{tabular}{c|c|c|c|c}
$g(x)$ & $\deg g$ & $k$ & $d_{\min}$ & Тип \\
\hline
$x+1$ & 1 & 6 & 2 & код чётности \\
$x^3+x+1$ & 3 & 4 & 3 & код Хэмминга \\
$x^3+x^2+1$ & 3 & 4 & 3 & код Хэмминга \\
$x^4+x^3+x^2+1$ & 4 & 3 & 4 & симплекс-код \\
$x^4+x^2+x+1$ & 4 & 3 & 4 & симплекс-код \\
$x^6+\dots+1$ & 6 & 1 & 7 & код повторения \\
\end{tabular}
\end{center}

Минимальные расстояния были определены перебором всех кодовых слов.

\section{Описание полученных кодов}

Рассмотренные конструкции дают все возможные двоичные циклические коды длины $7$:

\begin{itemize}
    \item Полином $x+1$ задаёт код, в котором сумма координат любого слова равна нулю (код чётности).
    
    \item Полиномы $x^3+x+1$ и $x^3+x^2+1$ приводят к двум эквивалентным реализациям кода Хэмминга $(7,4,3)$.
    
    \item Полиномы четвёртой степени определяют симплекс-коды $(7,3,4)$, являющиеся двойственными к кодам Хэмминга.
    
    \item Полином степени $6$ задаёт код повторения, содержащий только два слова:
    \[
    (0,0,\dots,0), \quad (1,1,\dots,1).
    \]
\end{itemize}

\section{Проверка разложения}

Факторизацию удобно проверить с помощью системы компьютерной алгебры.
Например, в Python с использованием библиотеки \texttt{sympy}:

\begin{verbatim}
import sympy as sp
x = sp.symbols('x')
p = sp.Poly(x**7 + 1, domain=sp.GF(2))
print(p.factor())
\end{verbatim}

\end{document}